\subsection{物理信息神经网络}
\label{sec:physics-informed}

物理信息神经网络(PINN)将物理定律作为损失函数中的约束嵌入，确保预测遵守已知的物理行为\cite{}。这种方法在结构健康监测中显示出前景，领域特定知识可指导学习过程。

降解趋势约束变分自编码器(DTC-VAE)代表了导波SHM物理约束深度学习的一个重要进展\cite{dtcvae2026}。通过在VAE潜在空间中引入单调性约束，DTC-VAE确保学习到的表示编码物理上有意义的降解趋势。
类似地，能量导向自编码器已被开发用于健康状态评估，利用能量守恒原理提高特征可解释性\cite{perry2026health}。

这些物理约束通过将领域特定不变量直接编码到网络架构或优化目标中来改善特征保留。
物理信息模型不仅依赖数据驱动学习，还能即使在噪声条件下也保持声学有意义的特征，如到达时间、振幅和波形形态。

尽管有这些进展，物理约束在激光超声去噪中仍很少被应用。现有的激光超声深度学习方法通常采用通用去噪目标，没有显式的物理正则化。
这一差距提供了一个机会：嵌入保留声学特征的基于物理的损失项可产生更好地支持缺陷检测的去噪信号，同时减少对大量信号平均的依赖。

% word count: ~190
